\section{Partiality and well-definedness}
\label{sec:barrel-wd}

This is the section in which the encoding differs most from the B method's own treatment of
partial operators.
\Cref{subsec:barrel-partial-operators} explains how definedness evidence is turned into a
proof argument, and what that buys; \Cref{subsec:barrel-own-wd} explains why \barrel{}
recomputes such evidence instead of importing the conditions Atelier~B can generate.

\subsection{Well-definedness as a proof argument}
\label{subsec:barrel-partial-operators}

Lean has no primitive notion of partiality: a declared function is total on its domain
type.
A partial B operator must therefore be turned into a total Lean function, and there are
essentially three ways to do so.
One may return an arbitrary value on arguments outside the intended domain; the resulting
theorems are then true of the Lean definition but not of the B operator it is meant to
denote, and nothing in the statement records the discrepancy.
One may restrict the argument type by a subtype, which changes the shape of every
surrounding expression.
Or one may keep the argument types unchanged and add a proof of definedness as an extra
argument.
\barrel{} takes the third option.

For each partial operator, a proposition named \leaninl{WD} states the evidence the operator
needs, and the operator takes an inhabitant of that proposition as an explicit argument.
It may then use classical choice to return the value whose existence the evidence
establishes.
\Cref{lst:barrel-partial-definitions} shows the pattern on minimum and relational
application. In both, the evidence is---up to unfolding---an existential statement, which is
what \leaninl{Classical.choose} consumes (\Cref{sec:lean}): \leaninl{isBoundedBelow} is one
as stated, and \leaninl{isInDomain} becomes one because \leaninl{dom f} is by definition
\leaninl!{x | ∃ y, (x, y) ∈ f}!, so \leaninl{x ∈ dom f} unfolds to
\leaninl{∃ y, (x, y) ∈ f}.

\begin{leanlisting}[title={Proof-relevant encodings of the minimum of a set and of
    relational application.}, label={lst:barrel-partial-definitions}, float=t]
structure min.WD {α} [PartialOrder α] (S : Set α) : Prop where
  isBoundedBelow : ∃ x ∈ S, ∀ y ∈ S, x ≤ y

noncomputable abbrev min {α} [PartialOrder α]
    (S : Set α) (wd : min.WD S) : α :=
  Classical.choose wd.isBoundedBelow

structure app.WD {α β} (f : SetRel α β) (x : α) : Prop where
  isPartialFunction : f ∈ tfun (dom f) (ran f)
  isInDomain        : x ∈ dom f

noncomputable abbrev app {α β}
    (f : SetRel α β) (x : α) (wd : app.WD f x) : β :=
  Classical.choose wd.isInDomain
\end{leanlisting}

The consequence is structural rather than methodological.
The Lean term denoting \(\appB f\,x\) has type \(\beta\) only when its third argument has
type \leaninl{app.WD f x}; there is no term \leaninl{app f x} to write down.
A well-definedness condition is thus not displayed \emph{next to} the main theorem: its
proof occurs \emph{inside} that theorem's statement.
The evidence may be obtained automatically, interactively, or by reusing an earlier
theorem, but it cannot be omitted, and it cannot be introduced later by a proof step that
forgot to check it.

\Cref{tab:barrel-wd-predicates} lists the five partial operators of the supported fragment
together with the evidence each of them requires.
Four consume that evidence through \leaninl{Classical.choose}; the field names are
historical and slightly understate their contents.
For the minimum, \leaninl{isBoundedBelow} asserts the existence of a least element, not
merely of a lower bound, which is what makes the chosen value unique.
For application, the second field alone already yields a value by choice; the first field,
which says that \(f\) is a total function \emph{on its own domain} and is how functionality
of a relation is phrased in this model, is what makes that value \emph{the} image of \(x\).
Cardinality is the exception, and it shows that the proof argument is not a mere guard:
from \leaninl{S.Finite} the definition builds a \leaninl{Fintype}
instance for \leaninl{S} and returns the cardinality of the resulting finite set, so the
proof genuinely participates in the elaboration of the value.

\begin{table}[t]
  \centering
  \small
  \begin{tabular}{lll}
    \toprule
    B operator & Lean operator & Well-definedness evidence \\
    \midrule
    \(\minB(S)\)      & \leaninl{min}  & \(S\) has a least element \\
    \(\maxB(S)\)      & \leaninl{max}  & \(S\) has a greatest element \\
    \(\cardB{S}\)     & \leaninl{card} & \(S\) is finite \\
    \(f(x)\)          & \leaninl{app}  & \(f \inB \domB(f) \tfunB \ranB(f)\) and \(x \inB \domB(f)\) \\
    \(\mathsf{size}(s)\) & \leaninl{size} & \(s\) is a sequence over its own range \\
    \bottomrule
  \end{tabular}
  \caption{Partial operators of the supported fragment and the well-definedness evidence
    they take as an argument.}
  \label{tab:barrel-wd-predicates}
\end{table}

\subsection{Why \barrel{} regenerates its own conditions}
\label{subsec:barrel-own-wd}

Atelier~B is perfectly capable of emitting well-definedness obligations, and does so on
request.
\barrel{} nonetheless recomputes them from the Lean term it is constructing, for three
reasons.

First, the evidence must be a \emph{subterm} of the main proposition, not a sibling of it.
The proof-obligation file records well-definedness goals as separate obligations, with no
term-level link between a particular occurrence of a partial operator and the side goal
meant to justify it; recovering that link would mean re-deriving, from the text of the
obligations, the association that the encoder is in the middle of establishing anyway.
Worse, the two sets of conditions are stated over different objects, since Atelier~B's
generator simplifies its side goals and drops hypotheses it judges irrelevant, so the
resulting formula is generally not the proposition that \barrel{}'s operator signature
demands.

Second, the imported conditions would be attached to the goal only.
An expression introduced during an interactive proof---a witness for an existential, an
instantiation of a universally quantified hypothesis---never passes through
proof-obligation generation again.
Any discipline based on precomputed side goals is therefore silent about the terms a user
writes. \barrel{} needs no second mechanism for them: by
\Cref{subsec:barrel-partial-operators}, a term containing a partial operator does not
elaborate until its definedness argument is supplied, wherever it is introduced.
\Cref{sec:barrel-atelier-comparison} examines this situation.

Third, reinstalling an imported condition would require the context in which the occurrence
was encountered.
Well-definedness obligations are emitted after the goals that need them, they are
identified only by their textual description, and a condition may depend on a preceding
conjunct of the very goal it belongs to.
The assertion
\[
  \existsB \var{G} \cdot \bigl(\var{G} \inB A \tfunB B \andB \var{G}(\var{x}) \inB B \bigr)
\]
produces the definedness condition
\[
  \forallB \var{G} \cdot \bigl(\var{G} \inB A \tfunB B \impB
    \var{x} \inB \domB(\var{G}) \andB \var{G} \inB \domB(\var{G}) \pfunB \ranB(\var{G}) \bigr),
\]
whose hypothesis is the first conjunct of the source assertion.
That context is available during encoding and nowhere else, and is precisely what the
reification mechanism of \Cref{subsec:barrel-reification} captures.

\begin{remark}
  Recomputing does not mean disagreeing.
  For integer extrema, Atelier~B phrases definedness as non-emptiness together with a
  finiteness condition on an intersection with the naturals or with their complement, while
  \barrel{} uses the order-theoretic condition that the set has a least or greatest
  element, which also applies to reals and to any partially ordered carrier.
  The library proves the two characterizations equivalent for the integer extrema---the
  only operators for which such a bridging theorem exists. For the minimum,
  \[
    \mathtt{min.WD}\ S
    \quad\iff\quad
    S \neq \varnothing
      \;\land\; S \cap (\mathtt{INTEGER} \mathbin{\backslash} \mathtt{NATURAL})
        \in \mathtt{FIN}\ \mathtt{INTEGER},
  \]
  and symmetrically for the maximum, with \(\mathtt{NATURAL}\) in place of its complement.
  What differs is not the mathematical content but where the condition is anchored.
\end{remark}
